\documentclass[11pt, a4paper]{report}

\usepackage[utf8]{inputenc}
\usepackage[T1]{fontenc}
\usepackage[french]{babel}

% Appel des packages dans le dossier preambule
\def\packagepath{../../../../preambule}   % path du package princial

\usepackage{\packagepath/page_layout} % <--- Nouveau
\usepackage{\packagepath/config_info}
\usepackage{\packagepath/cadres_cours}
\usepackage{\packagepath/zones_reponse}
\usepackage{\packagepath/config_ds}
\usepackage{\packagepath/config_maths}

% --- PILOTAGE ---
\def\visibleornot{invisible}    % visible or invisible


\def\documentpath{./Sources_Latex}
\graphicspath{{./Images/}}


\def\ptsexoA{0}



\begin{document}

\def\level{BTS}  % Classe à droite
\def\course{CIEL 1}
\def\chaptername{Nombres complexes} % Titre au centre
\def\docname{C112~--~TD 1}        % Identifiant à gauche

\pagestyle{DS_FP}

\NomPrenom{}

\bigskip

%%=============================================================
\exods{0} \medskip

\textbf{Compréhension et Forme Algébrique}: Soit un nombre complexe $z$ défini par:
\begin{center}
    $z = 3 - 4\i$
\end{center}

\begin{enumerate}
    \item \textbf{Terminologie}: Identifiez la partie réelle $\Ree(z)$ et la partie imaginaire $\Ima(z)$: \reponseLigne[4cm]{$\Ree(z) = 3$ et $\Ima(z) = -4$}
     \item \textbf{Conjugué}: Donnez l'écriture du nombre conjugué $\conj{z}$: \reponseLigne[4cm]{$\conj{z} = 3 + 4\i$} 
    \item \textbf{Module}: Calculez le module $|z|$: \reponseLigne[4cm]{$|z| = 5$}
    \item \textbf{Opposé}: Donnez l'affixe du point $P$ symétrique du point $M(z)$ par rapport à l'origine: \reponseLigne[4cm]{$z_P = -3 + 4\i$}
\end{enumerate}

\bigskip

%%=============================================================
\exods{0} \medskip

\textbf{Opérations et Puissances de $\i$}: 
Ecrire ces nombres complexes sous forme algébrique.

\begin{enumerate}
    \item $\i^2 = \reponseLigne[3cm]{$-1$}$
    \item $\i^3 = \reponseLigne[3cm]{$-\i$}$
    \item ${(1+\i)}^2 =  \reponseLigne[3cm]{$2\i$}$
    \item $z = \dfrac{3-\i}{4-3\i}=\reponseLigne[3cm]{$\dfrac{1}{5}+\dfrac{2}{5}\i$}$. 
\end{enumerate}


\bigskip

%%=============================================================
\exods{0} \medskip

\textbf{Résolution d'Équations}: 
On souhaite résoudre dans $\mathbb{C}$ l'équation du second degré suivante:
\begin{center}
    $4z^{2}-4z+5=0$
\end{center}

\begin{enumerate}
    \item Calculez le discriminant $\Delta = b^2 - 4ac$: \reponseLigne[3cm]{$\Delta = -64$}
    \item En déduire la nature des solutions (réelles ou complexes conjuguées)
    \begin{blocvisible}[height=1]
        $\Delta < 0$ donc les solutions sont complexes conjuguées.
    \end{blocvisible}

    \item Donnez les valeurs exactes des solutions $z_1$ et $z_2$:
    \begin{blocvisible}[height=1.5]
        \[z_1 = \frac{4 - \sqrt{-64}}{8} = \frac{4 - 8\i}{8} = \frac{1}{2} - \i \qquad z_2 = \frac{4 + \sqrt{-64}}{8} = \frac{4 + 8\i}{8} = \frac{1}{2} + \i\]
    \end{blocvisible}
\end{enumerate}

\bigskip
\pagebreak
\thispagestyle{DS_LP}
%%=============================================================
\exods{0} \medskip

\textbf{Formes Trigonométrique et Exponentielle}: 
Soit le nombre complexe $z_A = -2 + 2\i$. 

\begin{enumerate}
    \item Calculez le module $r = |z_A|$.
    
    \begin{blocvisible}[height=2]
        \[r = |z_A| = \sqrt{{(-2)}^2 + 2^2} = \sqrt{8} = 2\sqrt{2}\]
    \end{blocvisible}
    \item Déterminez un argument $\theta$ en utilisant les formules $\cos \theta = \dfrac{a}{r}$ et $\sin \theta = \dfrac{b}{r}$.
    \begin{blocvisible}[height=5]
        \[r = |z_A| = \sqrt{{(-2)}^2 + 2^2} = \sqrt{8} = 2\sqrt{2}\]
        \[\cos \theta = \frac{-2}{2\sqrt{2}} = -\frac{1}{\sqrt{2}} \qquad \sin \theta = \frac{2}{2\sqrt{2}} = \frac{1}{\sqrt{2}}\]
        \[\theta = \frac{\pi}{4} + \pi = \frac{5\pi}{4}\]       
    \end{blocvisible}
    \item Écrivez $z_A$ sous forme exponentielle ($r e^{i\theta}$).
    \begin{blocvisible}[height=2]
        \[z_A = |z_A| \cdot \e^{\i\theta}= 2\sqrt{2} \cdot e^{\dfrac{5\pi}{4}\i}\]
    \end{blocvisible}
    \item \textbf{Application Électronique}: Si une impédance complexe est $Z = 10 + 10i$, donnez sa forme exponentielle.
    \begin{blocvisible}[height=5]
        \[|Z| = \sqrt{10^2+10^2} = 10\sqrt{2}\]
        \[\cos \theta = \frac{10}{10\sqrt{2}}  = \frac{1}{\sqrt{2}} =\frac{\sqrt{2}}{2} \qquad \sin \theta = \frac{10}{10\sqrt{2}} = \frac{1}{\sqrt{2}}=\frac{\sqrt{2}}{2}\]
        \[\theta = \frac{\pi}{4}\]
        \[Z = 10\sqrt{2}\cdot e^{\i\dfrac{\pi}{4}}\]
    \end{blocvisible}   
\end{enumerate}

\end{document}